

\clearpage
\section{Metodología}
En esta sección se describe el dónde, cómo y con qué se va a desarrollar la investigación. De manera puntual, el dónde, hace referencia a la descripción del área física, territorio de estudio o localización de la investigación. El cómo, señala el procedimiento de la secuencia lógica de actividades para alcanzar cada uno de los objetivos específicos planteados. El con qué, describe el uso de insumos, materiales, equipos, necesarios para implementar la investigación. 
\subsection{Consideraciones éticas}
Aquí se debe redactar como se trataran los dato dentro del proyectos, si estos datos se refieren a personas se debe explicar como se hará la recolección para que no afecte las privacidad de los mismos, si no son datos de personas y son por ejemplos dataset de imágenes, sonidos, datos numéricos, etc, se debe describir como se recolectarán los datos o si es un dataset que esta en la nube, cuál es la licencia que tiene el mismo y si es permitido ocupar esos datos para la investigación que se pretende hacer.

\subsection{Resumen de la metodología}
En la Tabla \ref{tab:alcance} se presenta la metodología del proyecto, el mismo que contiene las actividades incluidas por cada objetivo, así como los entregables que se deben generar por cada unidad, los métodos y técnicas que se pretende utlizar y las herramientas y tecnologías con las que se desarrollará el proyecto.

\clearpage
\begin{landscape}
\begin{table}[h]
\centering
\caption{Metodología aplicada al proyecto}
\label{tab:alcance}
%\resizebox{\linewidth}{!}{%
\begin{tabular}{|c|p{5cm}|p{5cm}|p{5cm}|p{5cm}|}
\hline
\textbf{Objetivo} & \textbf{Actividades } & \textbf{Entregables} & \textbf{Métodos/Técnicas}& \textbf{Herramientas/Tecnologías} \\ \hline


\multirow{3}{*}{O.E.1} 
& Capturar imágenes en campo (en Loja) 
& Data set versión 1 
& Ingeniería de Datos: Obtención de los datos & Python, Colab\\ \cline{2-5}

& Etiquetar manualmente los datos basado en diagnóstico visual 
& Data set versión 2
& Ingeniería de Datos: Limpieza de los datos & Python \\ \cline{2-5}

& Preprocesar imágenes (resizing, normalización) 
& Data set verssión 3 
& Ingeniería de Datos: Obtención de los datos & Python\\ \hline

\multirow{2}{*}{O.E.2} 
& Ajustar un modelo CNN 
& Modelo ajustado versión 1. 
& Ingeniería de los modelos& Python, Colab \\ \cline{2-5}

& Evaluar el modelo con métricas cuantitativas 
& Modelo validado mediante validación cruzada con usuarios finales 
& Validación Cruzada& Python, Colab\\ \hline
\end{tabular}
%}
\end{table}
\end{landscape}


\clearpage


